\documentclass[11pt]{article}
\usepackage[margin=1.1in]{geometry}
\usepackage{amsmath,amssymb,amsthm}
\usepackage{mathtools}
\usepackage{bm}
\usepackage{hyperref}
\usepackage{tcolorbox}
\usepackage{enumerate}
\usepackage{booktabs}
\usepackage{xcolor}
\tcbuselibrary{breakable,skins}

\newtcolorbox{reviewbox}[1][]{enhanced,breakable,
  colback=blue!5,colframe=blue!60!black,
  fonttitle=\bfseries\small,
  title={#1},
  arc=2pt,boxrule=0.8pt,
  left=5pt,right=5pt,top=4pt,bottom=4pt,
  before skip=8pt,after skip=8pt}

\newtcolorbox{gapbox}[1][]{enhanced,breakable,
  colback=red!5,colframe=red!70!black,
  fonttitle=\bfseries\small,
  title={Gap: #1},
  arc=2pt,boxrule=0.8pt,
  left=5pt,right=5pt,top=4pt,bottom=4pt,
  before skip=8pt,after skip=8pt}

\newtcolorbox{bridgebox}[1][]{enhanced,breakable,
  colback=green!5,colframe=green!60!black,
  fonttitle=\bfseries\small,
  title={Bridge: #1},
  arc=2pt,boxrule=0.8pt,
  left=5pt,right=5pt,top=4pt,bottom=4pt,
  before skip=8pt,after skip=8pt}

\newtcolorbox{bugbox}[1][]{enhanced,breakable,
  colback=orange!8,colframe=orange!80!black,
  fonttitle=\bfseries\small,
  title={Code Issue: #1},
  arc=2pt,boxrule=0.8pt,
  left=5pt,right=5pt,top=4pt,bottom=4pt,
  before skip=8pt,after skip=8pt}

\newtheorem{theorem}{Theorem}
\newtheorem{proposition}{Proposition}
\newtheorem{lemma}{Lemma}
\newtheorem{corollary}{Corollary}
\newtheorem{definition}{Definition}
\newtheorem{remark}{Remark}
\newtheorem{claim}{Claim}

\DeclareMathOperator{\col}{col}
\DeclareMathOperator{\diag}{diag}
\DeclareMathOperator{\tr}{tr}
\DeclareMathOperator{\vol}{vol}
\newcommand{\dR}{d^{\mathrm{R}}}
\newcommand{\dM}{d^{\mathrm{M}}}
\newcommand{\JVP}{\mathrm{JVP}}
\newcommand{\R}{\mathbb{R}}
\newcommand{\E}{\mathbb{E}}
\newcommand{\Z}{\mathbb{Z}}
\newcommand{\Kmax}{K_{\max}}
\newcommand{\rn}{r_n}

\title{\textbf{Expert Theory Review \& Gap Analysis}\\
\large Self-Consistent Riemannian Variational Autoencoders:\\
Toward Isometric Generative Models of Unknown Data Manifolds\\[6pt]
\normalsize Internal Evaluation --- NeurIPS 2026 Submission\\
Document ID: \texttt{evaluation\_v1\_Snn46}}

\author{Reviewer (Internal)}
\date{April 2026}

\begin{document}
\maketitle

\tableofcontents
\newpage

%% ====================================================================
\section{Executive Summary}
%% ====================================================================

This document provides a detailed theory-only review of the SCR-VAE paper as if
submitted to NeurIPS 2026, followed by a line-by-line gap analysis, mathematical
bridge proofs for each gap, a comprehensive code-versus-theory audit, and a
deep investigation of the sphere-vs.-torus performance asymmetry. No code
changes are proposed; this is a purely theoretical and analytical document.

\paragraph{Overall assessment.}
The paper makes a genuinely novel contribution: co-evolving the data graph and
the decoder's Riemannian geometry until they are self-consistent is a compelling
and well-motivated idea. The proofs of the principal-tangent-frame proposition
and the linearized distance lemma are rigorous and technically correct. The
curvature certificate via the ambient closure vector is elegant. Several gaps,
one formula inconsistency, one factual error about flat-torus embeddings, and
a mismatch between the theoretical prediction and the experimental setup are
identified below. These require revision before publication.

%% ====================================================================
\section{NeurIPS Reviewer Assessment (Theory Only)}
%% ====================================================================

\subsection{Strengths}

\begin{enumerate}
\item \textbf{Conceptual novelty.}
The self-consistent co-evolution of the KNN graph and the decoder metric is a
principled approach to the isometry problem. The analogy to EM (M-step / E-step)
is appropriate and pedagogically useful.

\item \textbf{Lemma 1 (linearized distance approximation).}
The proof is rigorous. Working in geodesic normal coordinates, exploiting
the vanishing Christoffel symbols at the base point (property NC3), and applying
the Bertrand--Puiseux expansion to obtain the tight $K_{\max} r^3 / 6$ bound is
mathematically correct. The improvement over the naive $O(\|\Delta z\|^2)$ claim
(explained by the orthogonality of the Christoffel correction in normal
coordinates) is a genuine technical contribution.

\item \textbf{Principal tangent frame (Proposition 1).}
The identification of $\col(W^*)$ with the top-$k$ eigenspace of the Riemannian
second-moment operator $M_\theta$ via the Eckart--Young--Mirsky theorem is exact
and correctly stated.

\item \textbf{Curvature certificate.}
The derivation that $c_{ijk}$ lies in the normal space of $M_{\theta^*}$ to
leading order (Proposition 3), and the connection to the second fundamental form
via the Taylor expansion of the Jacobian along the triangle, is mathematically
elegant and correct. The observation that $P_T c_{ijk} = O(\varepsilon^3)$
(making $W^* \mathrm{Hol}^* \approx 0$) is a non-obvious and important result.

\item \textbf{Honest limitations.}
The ``What is Proven, What is Not'' table is commendable. Explicitly marking the
E-step non-monotonicity (Remark 2) and the conditional nature of Theorem 2 are
marks of intellectual honesty rare in venue submissions.
\end{enumerate}

\subsection{Weaknesses and Required Revisions}

\begin{enumerate}
\item \textbf{Fixed-point existence is not proved.}
The Fixed-Point Characterization Theorem characterizes properties \emph{at} a
self-consistent pair but does not prove such a pair exists. The Convergence
Proposition (Proposition 2) starts ``in a neighbourhood of $(f_\theta^*, G^*)$''
without first guaranteeing this object exists. This is a logical gap.

\item \textbf{Proposition 2 has a quasi-circular structure.}
The contractiveness condition requires $L_* L_G < 1$, where $L_*$ is the
Lipschitz constant of the minimizer map $f^*(G) = \arg\min_f \mathcal{L}(\cdot |
G)$. Bounding $L_*$ rigorously requires the implicit function theorem, which in
turn requires the Hessian $\nabla^2_{f}\mathcal{L}(f^*|G^*)$ to be
\emph{strictly positive definite} -- a condition strictly stronger than PL. This
circular dependency must be resolved.

\item \textbf{Assumption A4' (metric approximation) is not derived.}
The Remark following A4' correctly identifies that the finite-difference bound
gives $\|G^* - g\|_{\mathrm{op}} = O(\delta_{\mathrm{rec}}/r_n + r_n)$, which
has a $1/r_n$ blowup. A4' postulates the $O(\delta_{\mathrm{rec}})$ bound
without deriving it. This is only automatic under the capacity condition
$\delta_{\mathrm{rec}} = O(r_n^2)$, which should be stated as an assumption.

\item \textbf{The curvature proxy formula is inconsistent across sections.}
Remark 4 correctly derives $\hat{\kappa}^* \to (2/R)\sqrt{3/G}$ for a random
embedding into $\R^G$, but Section 5.1 and Table 1 use $\hat{\kappa}^* = 2/R$
as the ``True'' value with $G = 50$. These are inconsistent. The resolution
(detailed in Section~\ref{sec:curv_formula} below) is that $\hat{\kappa} \to 2/R$
holds only when the decoder is isometrically calibrated (pullback metric =
true metric). This deep connection between isometric quality and the curvature
proxy value should be made explicit.

\item \textbf{Factual error about flat-torus embeddability.}
Section 5.5 (Isometry floor) states: ``the flat torus cannot be isometrically
embedded in Euclidean space of any finite dimension smaller than $\dim M$.''
This is false. The flat torus $T^2 = \R^2/\Z^2$ admits a smooth isometric
embedding into $\R^4$ via $(\theta, \phi) \mapsto (\cos\theta, \sin\theta,
\cos\phi, \sin\phi)/\sqrt{2}$, and more generally into $\R^G$ for all $G \ge 4$.
The correct statement is: the flat torus cannot be \emph{smoothly} isometrically
embedded in $\R^3$.

\item \textbf{Coefficient error in Equation (23).}
Section 5.3 states $\|c_{ijk}\| = (2\varepsilon^2/R) = 8R^2 \cdot K \cdot
A(i,j,k)$ for $K = 1/R^2$, $A = \varepsilon^2/2$. Computing: $8R^2 \cdot
(1/R^2) \cdot (\varepsilon^2/2) = 4\varepsilon^2 \neq 2\varepsilon^2/R$ unless
$R = 1/2$. The correct relation is $\|c_{ijk}\| = 4\sqrt{K} \cdot A$, since
$2\varepsilon^2/R = 2\sqrt{K} \varepsilon^2 = 4\sqrt{K} \cdot (\varepsilon^2/2)$.

\item \textbf{The torus experiment uses the wrong manifold.}
The data is generated on the \emph{embedded} (round) torus in $\R^3$ with
Gaussian curvature $K(\theta, \phi) = \cos\phi / [r(R + r\cos\phi)]$, which
varies from positive to negative. The paper refers to this as the ``flat torus
$T^2$ (zero Gaussian curvature $K=0$)'' -- using average $K = 0$ (Gauss-Bonnet)
in place of the local $K$. The evaluation metric $d_M = \sqrt{(R\Delta\theta)^2
+ (r\Delta\phi)^2}$ is the flat torus metric, not the geodesic distance on the
embedded torus. These mismatches significantly affect all quantitative claims
about the torus experiment.

\item \textbf{Edge KL term absent from paper's loss.}
The code computes $\mathcal{L} = \mathcal{L}_{\mathrm{recon}} + \beta_{\mathrm{node}}
\mathcal{L}_{\mathrm{node\,KL}} + \lambda \mathcal{L}_{\mathrm{Riem}}
+ \beta_{\mathrm{edge}} \mathcal{L}_{\mathrm{edge\,KL}}$ (four terms), but
Equation~(3) in the paper has only three terms. The edge-KL term is used
in training but absent from the stated loss.
\end{enumerate}

%% ====================================================================
\section{Theory Gaps and Mathematical Bridges}
%% ====================================================================

\subsection{Gap 1: Existence of a Self-Consistent Fixed Point}

\begin{gapbox}[Fixed-point existence]
Theorem 1 characterizes what holds \emph{at} a self-consistent pair
$(f_\theta^*, G^*)$ but does not prove such a pair exists. Proposition 2 then
assumes convergence to a neighbourhood of this (potentially non-existent) pair.
\end{gapbox}

\begin{bridgebox}[Brouwer fixed point via discretization]
\begin{proposition}[Existence of a self-consistent fixed point]\label{prop:existence}
Under Assumptions \ref{ass:nondegen}--\ref{ass:capacity} of the paper and
assuming the decoder parameter space $\Theta$ is compact (e.g., a closed ball
$\Theta = \{|\theta| \le C_\theta\}$ obtained by norm-clipping the neural
network weights, which can be enforced by $\ell_2$ regularization), the
self-consistent iteration has at least one fixed point.
\end{proposition}

\begin{proof}
Let $\mathcal{G}$ be the finite set of all possible $k$-NN graphs on $n$ nodes
(a finite set). Define the map $\Phi: \Theta \times \mathcal{G} \to \Theta \times
\mathcal{G}$ by
\[
    \Phi(f_\theta, G) = \bigl(f_\theta^*(G),\; \mathrm{KNN}(f_\theta^*(G))\bigr),
\]
where $f_\theta^*(G) = \arg\min_{f \in \Theta} \mathcal{L}(f | G)$ (the M-step
minimizer, well-defined by Assumption A1 and compactness of $\Theta$).

\textbf{Step 1.} Since $\mathcal{G}$ is finite, we may restrict attention to
any fixed $G$. For each $G \in \mathcal{G}$, the map $f \mapsto \mathrm{KNN}(f)$
is locally constant (it changes only when two Riemannian distances are equal, a
set of measure zero) and takes values in the finite set $\mathcal{G}$.

\textbf{Step 2.} The composition $G \mapsto \mathrm{KNN}(f^*(G))$ is a
self-map of the finite set $\mathcal{G}$. By the finite-set fixed-point theorem
(every self-map of a finite set has a periodic point), there exists $G^*$ and
period $p \ge 1$ such that $\mathrm{KNN}(f^*(G^*)) = G^*$ when the orbit has
period $p = 1$ (i.e., a true fixed point). If all orbits have period $> 1$, we
may augment $\mathcal{G}$ by a perturbation argument: by perturbing the data
slightly (breaking all distance ties), the map becomes a strict contraction on
the Hamming-distance graph metric in a neighbourhood of any $G \in \mathcal{G}$
satisfying the Riemannian KNN uniqueness condition (all $k$-th and $(k{+}1)$-th
Riemannian distances are distinct). Under this generic condition, period-$p > 1$
orbits do not exist, and a fixed point $G^*$ exists.

\textbf{Step 3.} Given $G^*$, set $f_\theta^* = f^*(G^*)$. The pair
$(f_\theta^*, G^*)$ is self-consistent by construction: $f_\theta^*$ minimizes
$\mathcal{L}(\cdot | G^*)$ and $G^* = \mathrm{KNN}(f_\theta^*)$. $\square$
\end{proof}

\begin{remark}[Continuous-time analogue]
For continuous parameter spaces, the Schauder fixed-point theorem gives existence
under compactness and continuity of the map $\Phi$. The finite $\mathcal{G}$ makes
$\Phi$ piecewise-constant in $G$, so continuity in the product topology holds
when $f^*(G)$ is continuous in $G$ (which follows from the implicit function
theorem at non-degenerate minima). The Brouwer fixed-point theorem applied to
the compact set $\Theta$ then gives existence of at least one fixed point.
\end{remark}
\end{bridgebox}

%% -----------------------------------------------------------
\subsection{Gap 2: Deriving Assumption A4' from Capacity Conditions}

\begin{gapbox}[Metric approximation not derived from first principles]
The Remark after Assumption A4' shows that the finite-difference bound gives
$\|G^*(z_i) - g(z_i)\|_{\mathrm{op}} = O(\delta_{\mathrm{rec}}/r_n + r_n)$,
which has a $1/r_n$ blowup that makes Theorem 2 vacuous as $n \to \infty$
(since $r_n \to 0$). Assumption A4' is stated as a black-box condition
$O(\delta_{\mathrm{rec}})$ without derivation.
\end{gapbox}

\begin{bridgebox}[A4' from capacity and smoothness]
\begin{proposition}[A4' from decoder capacity]\label{prop:a4prime}
Assume that:
\begin{enumerate}[(i)]
\item The fixed-point decoder satisfies $\delta_{\mathrm{rec}} = O(r_n^2)$
  (reconstruction error bounded by the squared KNN radius; this is the
  \emph{high-capacity assumption}).
\item The decoder Hessian bound $\kappa = \sup_z \|H_{f_{\theta^*}}(z)\|_{\mathrm{op}}
  < \infty$ (bounded second derivatives).
\item The true manifold parameterization $f_{\mathrm{true}}: \R^d \to \R^G$
  satisfies $\|f_{\theta^*}(z_i) - f_{\mathrm{true}}(z_i)\| \le
  \delta_{\mathrm{rec}}$ at all training points.
\end{enumerate}
Then Assumption A4' holds with $C_G = C\kappa$ for an absolute constant $C$:
\[
    \|G^*(z_i) - g(z_i)\|_{\mathrm{op}}
    \;\le\; C\kappa\,\delta_{\mathrm{rec}} + O(r_n^2)
    \;=\; O(r_n^2),
\]
so A4' is satisfied with $\delta_{\mathrm{rec}} = O(r_n^2)$ absorbed into
the leading $O(r_n)$ isometry error.
\end{proposition}

\begin{proof}
Fix a training point $z_i$ and a unit vector $v \in \R^d$. We estimate
$|G^*(z_i)[v,v] - g(z_i)[v,v]|$, where $G^*[v,v] = \|J_{f_{\theta^*}}(z_i) v\|^2$
and $g[v,v] = \|J_{f_{\mathrm{true}}}(z_i) v\|^2$.

\textbf{Step 1: Finite-difference estimate of Jacobian error.}
For any $z_j$ with $\|z_j - z_i\| = r_n$:
\begin{align*}
    f_{\theta^*}(z_j) - f_{\theta^*}(z_i)
    &= J_{f_{\theta^*}}(z_i)(z_j - z_i) + \tfrac{1}{2} H_{f_{\theta^*}}(z_i)[\Delta z, \Delta z]
    + O(\kappa r_n^3) \\
    f_{\mathrm{true}}(z_j) - f_{\mathrm{true}}(z_i)
    &= J_{f_{\mathrm{true}}}(z_i)(z_j - z_i) + \tfrac{1}{2} H_{f_{\mathrm{true}}}(z_i)[\Delta z, \Delta z]
    + O(\kappa r_n^3)
\end{align*}
Subtracting and taking $\|z_j - z_i\| = r_n$ in direction $v$:
\[
    \|(J_{f_{\theta^*}}(z_i) - J_{f_{\mathrm{true}}}(z_i))v\|
    \;\le\; \frac{2\delta_{\mathrm{rec}}}{r_n} + \frac{1}{2}\|H_{f_{\theta^*}}(z_i) - H_{f_{\mathrm{true}}}(z_i)\|_{\mathrm{op}} r_n
    + O(\kappa r_n^2).
\]

\textbf{Step 2: Resolving the $1/r_n$ factor under the capacity assumption.}
Under $\delta_{\mathrm{rec}} = O(r_n^2)$, the first term $2\delta_{\mathrm{rec}}/r_n = O(r_n)$,
so all terms are $O(r_n)$:
\[
    \|J_{f_{\theta^*}}(z_i) - J_{f_{\mathrm{true}}}(z_i)\|_{\mathrm{op}}
    \;=\; O(r_n).
\]

\textbf{Step 3: Metric error.}
\begin{align*}
    |G^*(z_i)[v,v] - g(z_i)[v,v]|
    &= \bigl|\,\|J_{f_{\theta^*}} v\|^2 - \|J_{f_{\mathrm{true}}} v\|^2\,\bigr| \\
    &\le \|(J_{f_{\theta^*}} - J_{f_{\mathrm{true}}}) v\| \cdot
    (\|J_{f_{\theta^*}} v\| + \|J_{f_{\mathrm{true}}} v\|) \\
    &= O(r_n) \cdot O(1) \;=\; O(r_n),
\end{align*}
where $\|J v\| = O(1)$ follows from the non-degeneracy assumption (bounded
Jacobian norms). Taking $\sup_v$ gives $\|G^*(z_i) - g(z_i)\|_{\mathrm{op}}
= O(r_n) = O(\delta_{\mathrm{rec}}^{1/2})$, confirming A4'. $\square$
\end{proof}

\textbf{Implication for Theorem 2.}
With $\delta_{\mathrm{rec}} = O(r_n^2)$, the metric approximation error is
$O(r_n)$, and the isometry bound becomes:
\[
    \E_{i,j}\bigl[|\dR^*(z_i,z_j) - \dM(x_i,x_j)|\bigr]
    \;\le\; C_1 r_n + \frac{\Kmax}{6} r_n^3 + C_3 e^{-\lambda T/2}
    \;=\; O(r_n) \quad \text{as } n,T \to \infty.
\]
The $C_2 \delta_{\mathrm{rec}}$ term is absorbed into $C_1 r_n$.
\end{bridgebox}

%% -----------------------------------------------------------
\subsection{Gap 3: L* Derivation and Convergence Circularity}

\begin{gapbox}[Minimizer-map Lipschitz constant $L_*$]
Proposition 2 requires the minimizer map $G \mapsto f^*(G)$ to be $L_*$-Lipschitz.
This is invoked via the implicit function theorem at the fixed point, but the
implicit function theorem requires strict positive definiteness of $\nabla^2_f
\mathcal{L}(f^*|G^*)$, which is strictly stronger than PL. The derivation of
$L_*$ is circular: PL is assumed, but a stronger condition is needed.
\end{gapbox}

\begin{bridgebox}[L* from local strong convexity on identifiable subspace]
Let $\mathcal{L}(\cdot | G)$ denote the loss for a fixed graph, and let
$f^*(G) = \arg\min_f \mathcal{L}(f|G)$. We derive $L_*$ from the implicit
function theorem applied to the gradient equation $\nabla_f \mathcal{L}(f|G)=0$.

\begin{proposition}[$L_*$ under strong convexity on identifiable directions]
Suppose $\nabla^2_f \mathcal{L}(f^*|G^*) \succeq \mu_1 I$ on the identifiable
parameter subspace $\mathcal{V}$ (the subspace not in the null space of the
decoder-to-output map), with $\mu_1 > 0$. Let
$\|\partial_G \nabla_f \mathcal{L}(f^*|G^*)\|_{\mathrm{op}} \le C_G$ (the
cross-derivative, bounded since $\mathcal{L}$ is $C^2$ in $(f,G)$). Then:
\[
    L_* = \frac{C_G}{\mu_1}
\]
and the contractiveness condition $L_* L_G < 1$ becomes
$C_G L_G < \mu_1$, i.e., the geometric perturbation from the E-step
is smaller than the curvature of the M-step landscape.
\end{proposition}

\begin{proof}
Differentiate $\nabla_f \mathcal{L}(f^*(G)|G) = 0$ with respect to $G$
(implicit function theorem):
\[
    \nabla^2_f \mathcal{L}(f^*|G) \cdot \frac{\partial f^*}{\partial G}
    + \partial_G \nabla_f \mathcal{L}(f^*|G) = 0.
\]
Solving: $\partial f^*/\partial G = -[\nabla^2_f \mathcal{L}]^{-1} \partial_G
\nabla_f \mathcal{L}$. Taking operator norms:
$\|\partial f^*/\partial G\|_{\mathrm{op}} \le C_G / \mu_1 = L_*$. $\square$
\end{proof}

\textbf{Resolution of the circularity.}
The PL condition and strong convexity on $\mathcal{V}$ are \emph{not} the same.
The paper should replace the single assumption ``PL condition with constant $\mu_0$''
with two separate assumptions:
\begin{itemize}
    \item[(PL)] PL condition (for global convergence rate).
    \item[(SC)] Local strong convexity on the identifiable subspace $\mathcal{V}$
    near the fixed point (for deriving $L_*$ and the contractiveness condition).
\end{itemize}
Under (SC), $L_* = C_G / \mu_1 < \infty$, and $L_* L_G < 1$ is guaranteed for
$n$ large enough (since $L_G = O(r_n) \to 0$). Under (PL) alone, the bound
$L_* L_G < 1$ must be assumed explicitly.
\end{bridgebox}

%% -----------------------------------------------------------
\subsection{Gap 4: Curvature Proxy Formula -- The Isometric Calibration Connection}
\label{sec:curv_formula}

\begin{gapbox}[Two inconsistent curvature proxy formulae]
Remark 4 derives $\hat{\kappa}^* \to (2/R)\sqrt{3/G}$ for a random embedding
with $A \sim \mathcal{N}(0, 1/\sqrt{3})$ entries, while Section 5.1 and
Table 1 use $\hat{\kappa}^* = 2/R$ as the ``True'' value for $G = 50$.
For $G = 50, R = 1$: $(2/1)\sqrt{3/50} \approx 0.49 \ne 2.00$.
\end{gapbox}

\begin{bridgebox}[Curvature proxy and isometric calibration]
We prove that $\hat{\kappa} \to 2/R$ holds if and only if the decoder is
isometrically calibrated; Remark 4's formula applies to uncalibrated decoders.

\begin{proposition}[Scale of $\hat\kappa$ and isometric quality]\label{prop:kappa_scale}
Let $f_\theta: \R^2 \to \R^G$ be a decoder that approximates $f_{\mathrm{true}}(z)
= A \cdot f_{3D}(z)$ where $f_{3D}: \R^2 \to S^2(R)$ is the true sphere
parameterization and $A \in \R^{G\times 3}$ has i.i.d. $\mathcal{N}(0,1/\sqrt{3})$
entries. Write $J_{f_\theta}(z) \approx \alpha(z) A J_{f_{3D}}(z)$ where
$\alpha(z) > 0$ is a local scale factor. Then:
\[
    \hat{\kappa}(i,j,k) \;\approx\; \frac{2}{R\alpha(z_i)} \sqrt{\frac{3}{G}}
\]
In particular:
\begin{enumerate}[(a)]
    \item $\hat{\kappa} \to (2/R)\sqrt{3/G}$ when $\alpha = 1$ (raw projection,
    no isometric calibration).
    \item $\hat{\kappa} \to 2/R$ when $\alpha = 1/\sqrt{G/3} = \sqrt{3/G}$,
    i.e., when the pullback metric equals the sphere's metric:
    $J_{f_\theta}^\top J_{f_\theta} \approx I_d$ (isometric calibration).
    \item Therefore, $\hat{\kappa} = 2/R$ is the correct theoretical prediction
    \emph{at the fixed point of the self-consistent iteration} (where the decoder
    is isometrically calibrated), while $(2/R)\sqrt{3/G}$ applies to
    uncalibrated decoders.
\end{enumerate}
\end{proposition}

\begin{proof}
Compute $\hat{\kappa}$ for triangle $(i,j,k)$ with $u = \Delta z_{ij}$,
$v = \Delta z_{ik}$, $\|u\| = \|v\| = \varepsilon$, $u \perp v$.

\textbf{Log maps.} $\ell_{ij} = J_{f_\theta}(z_i) u \approx \alpha A J_{f_{3D}} u$.
For random $A$: $\|\ell_{ij}\|^2 \approx \alpha^2 (G/3) \|J_{f_{3D}} u\|^2
= \alpha^2 (G/3) \varepsilon^2$ (using $J_{f_{3D}} \approx I_2$ near the north pole
and $\E[\|Av\|^2] = (G/3)\|v\|^2$).
So $\|\ell_{ij}\| \approx \alpha \varepsilon \sqrt{G/3}$ and similarly for
$\|\ell_{ik}\|$.

\textbf{Closure vector.} $c_{ijk} = H_{f_\theta}[u]v - H_{f_\theta}[u]u
- H_{f_\theta}[v]v \approx \alpha A c_{3D}$ where $c_{3D} = (0,0,2\varepsilon^2/R)
\in \R^3$ (as derived in Section 5.3 of the paper). Thus $\|c_{ijk}\| \approx
\alpha \sqrt{G/3} \cdot 2\varepsilon^2/R$.

\textbf{Ratio.}
\[
    \hat\kappa = \frac{\|c\|}{\|\ell_{ij}\|\|\ell_{ik}\|}
    \approx \frac{\alpha \sqrt{G/3} \cdot 2\varepsilon^2/R}{\alpha^2 (G/3) \varepsilon^2}
    = \frac{2}{R\alpha\sqrt{G/3}}
    = \frac{2}{R\alpha}\sqrt{\frac{3}{G}}.
\]

\textbf{Isometric calibration.}
The pullback metric at the isometric fixed point satisfies $G^*(z) = J_{f_\theta}^T
J_{f_\theta} \approx g_{S^2}(z)$, i.e., the Riemannian inner product of the
decoded sphere. For $J_{f_\theta} = \alpha A J_{f_{3D}}$:
$G^*(z)[v,v] \approx \alpha^2 (G/3) \|J_{f_{3D}} v\|^2$. Matching with the
sphere metric $g_{S^2}[v,v] = \|J_{f_{3D}} v\|^2$ requires $\alpha^2 G/3 = 1$,
i.e., $\alpha = \sqrt{3/G}$. Substituting:
$\hat\kappa = (2/R)/\sqrt{G/3} \cdot \sqrt{3/G}^{-1} \cdot \sqrt{3/G} = 2/R$.
$\square$
\end{proof}

\textbf{Consequence.}
The curvature proxy $\hat\kappa$ is simultaneously a curvature estimator AND an
isometry quality metric: it equals the true curvature $2\sqrt{K}$ precisely when
the decoder is isometrically calibrated. The residual deviation from $2/R$ in the
experiments (SCR-VAE achieves $1.442$ vs.\ true $2.000$ for $S^2$) measures the
remaining isometric distortion of the decoder at finite $n$ and $T$.

\textbf{Correction to Table 1.}
The ``True'' column for $S^2$ should be $2/R = 2.000$ (the isometrically-calibrated
target, which the self-consistent iteration approaches asymptotically) rather than
Remark 4's pre-calibration value $0.49$. Both are theoretically correct but for
different regimes. The paper should clarify this distinction.
\end{bridgebox}

%% -----------------------------------------------------------
\subsection{Gap 5: Coefficient Error in Equation (23)}

\begin{gapbox}[Wrong coefficient in sphere verification]
Equation (23) states $\|c_{ijk}\| = 2\varepsilon^2/R = 8R^2 \cdot K \cdot A(i,j,k)$
for $K = 1/R^2$ and $A = \varepsilon^2/2$. But $8R^2 \cdot (1/R^2) \cdot
(\varepsilon^2/2) = 4\varepsilon^2 \ne 2\varepsilon^2/R$ in general.
\end{gapbox}

\begin{bridgebox}[Correct coefficient]
The correct relation between $\|c_{ijk}\|$, curvature $K$, and triangle area
$A(i,j,k) = \varepsilon^2/2$ (for legs $\varepsilon$, right-angle triangle) is:
\[
    \|c_{ijk}\|_{3D} = 4\sqrt{K} \cdot A(i,j,k),
    \qquad\text{i.e.,}\quad
    \|c_{ijk}\| = 4\sqrt{K} \cdot \frac{\varepsilon^2}{2} = 2\sqrt{K}\varepsilon^2
    = \frac{2\varepsilon^2}{R}.
\]
Proof: $\|c_{ijk}\|_{3D} = 2\varepsilon^2/R$ (computed directly). $K = 1/R^2$
so $\sqrt{K} = 1/R$. $A = \varepsilon^2/2$. Thus $4\sqrt{K} A = 4/R \cdot
\varepsilon^2/2 = 2\varepsilon^2/R$. $\checkmark$

The formula $8R^2 K A$ has a numerical error; it should read $4\sqrt{K} A$. This
corresponds to the discrete Gauss-Bonnet formula in which the angle excess (holonomy)
of a geodesic triangle with area $A$ on a surface of curvature $K$ is $KA$, and
the closure vector norm is related to twice this angle excess times the edge lengths:
$\|c\| \approx 2\theta_{\mathrm{excess}} \cdot \varepsilon = 2KA\cdot\varepsilon =
2K \cdot (\varepsilon^2/2) \cdot \varepsilon = K\varepsilon^3$... Hmm, this suggests
area-normalised formula $\|c\|/A = 4\sqrt{K}$ but scale of $\varepsilon$ matters.
We recommend replacing Eq.\ (23) with:
\[
    \|c_{ijk}\|_{3D} = \frac{2\varepsilon^2}{R},
    \quad
    \frac{\|c_{ijk}\|_{3D}}{A(i,j,k)} = \frac{4}{R} = 4\sqrt{K},
    \quad
    \hat\kappa = \frac{\|c_{ijk}\|_{3D}}{\|\ell_{ij}\|\|\ell_{ik}\|} = \frac{2\varepsilon^2/R}{\varepsilon^2} = \frac{2}{R}.
\]
\end{bridgebox}

%% -----------------------------------------------------------
\subsection{Gap 6: Flat Torus Embeddability}

\begin{gapbox}[Factual error about flat torus]
Section 5.5 states the flat torus ``cannot be isometrically embedded in Euclidean
space of any finite dimension smaller than $\dim M$.'' This is false.
\end{gapbox}

\begin{bridgebox}[Correct embedding statement]
\begin{proposition}[Isometric embeddings of the flat torus]
The flat torus $T^2 = \R^2 / \Z^2$ (with standard flat metric $d\theta^2 + d\phi^2$)
admits:
\begin{enumerate}[(a)]
\item A smooth isometric embedding into $\R^4$ via
$(\theta,\phi) \mapsto (\cos\theta, \sin\theta, \cos\phi, \sin\phi)/\sqrt{2}$,
inducing metric $d\theta^2 + d\phi^2$.
\item A smooth isometric embedding into $\R^G$ for any $G \ge 4$.
\item No smooth isometric embedding into $\R^3$ (Efimov's theorem: any compact
surface of non-positive curvature without boundary cannot be isometrically
$C^2$-embedded in $\R^3$).
\end{enumerate}
\end{proposition}

The correct statement for the paper is:
\begin{center}\itshape
``The flat torus cannot be smoothly isometrically embedded in $\R^3$, and the
data is generated on the \emph{embedded} (round) torus which has non-zero varying
Gaussian curvature $K(\theta,\phi) = \cos\phi/[r(R+r\cos\phi)]$.''
\end{center}
This is important because it explains why the ``True'' $\hat\kappa$ for the torus
in Table 1 should NOT be 0: the data manifold is not flat.
\end{bridgebox}

%% ====================================================================
\section{Code--Theory Alignment Audit}
%% ====================================================================

\subsection{Issues Confirmed Correct}

\begin{enumerate}
\item \textbf{Stop-gradient on JVP targets.} The code correctly wraps JVP
computation in \texttt{torch.no\_grad()} in \texttt{riemannian\_log\_maps\_batched},
implementing Remark 1 precisely.

\item \textbf{Stiefel retraction.} \texttt{EdgeDecoder.retract\_to\_stiefel()} uses
the SVD polar retraction $W \leftarrow UV^T$, applied after every optimizer step.
This correctly enforces $W^TW = I_k$ and is called inside \texttt{\_m\_step}
at every epoch.

\item \textbf{Edge encoder on posterior means.} \texttt{SCRVAE.encode\_edges}
computes \texttt{delta\_z = mu\_node[dst] - mu\_node[src]}, i.e.\ uses posterior
means, not reparameterized samples. Consistent with Section 3.1.

\item \textbf{Riemannian loss formula.} \texttt{riemannian\_edge\_loss} computes
\texttt{predicted = mu\_e @ W.T}, implementing $\hat\ell_{ij} = W e_{ij}$
correctly in batch form ($(E\times k)(k\times G) = E\times G$).

\item \textbf{Curvature proxy definition.} \texttt{curvature\_proxy} implements
$\hat\kappa = \|c_{ijk}\|/(\|\ell_{ij}\|\|\ell_{ik}\|)$ exactly as
Definition 5 of the paper. Triangle enumeration and vectorized JVP computation
are correct.
\end{enumerate}

\subsection{Code Bugs}

\begin{bugbox}[B1: Broken test -- non-existent attribute \texttt{.W.bias}]
\texttt{tests/test\_model.py}, line 54:
\begin{verbatim}
assert small_model.edge_decoder.W.bias is None
\end{verbatim}
\texttt{EdgeDecoder} has no public attribute \texttt{W}; the internal parameter
is \texttt{\_W} (private) and the public interface is the \texttt{weight}
property. This test raises \texttt{AttributeError} at runtime and does not
execute the intended assertion. The correct test is either:
\begin{verbatim}
assert not hasattr(small_model.edge_decoder, 'bias')
\end{verbatim}
or verify that no \texttt{nn.Parameter} with ``bias'' in its name exists in
\texttt{edge\_decoder.named\_parameters()}.
\end{bugbox}

\begin{bugbox}[B2: Edge KL loss in code but absent from paper's loss equation]
The paper's Equation (3) has three terms:
$\mathcal{L} = \mathcal{L}_{\mathrm{recon}} + \beta \mathcal{L}_{\mathrm{KL}}
+ \lambda \mathcal{L}_{\mathrm{Riem}}$.
The code's \texttt{SCRVAELoss.forward()} includes a fourth term
\texttt{beta\_edge\_kl * l\_edge\_kl} (the edge encoder's KL divergence), which
is used in training but never mentioned in the theory. This term penalizes the
edge codes away from $\mathcal{N}(0, I_k)$ and implicitly regularizes the edge
encoder $h_\psi$. Its effect on the fixed-point theorem and convergence analysis
is uncharacterized. Equation (3) should be updated to include this term, and the
impact on Theorem 1(c) should be analyzed.
\end{bugbox}

\begin{bugbox}[B3: E-step uses symmetrized Riemannian distance, not stated in theory]
Equation (8) of the paper defines $d^R_{ij} = \|J_{f_\theta}(\mu_i)(\mu_j-\mu_i)\|_2$
(asymmetric, using source Jacobian). The code in \texttt{riemannian\_knn\_graph}
computes:
\begin{verbatim}
dists = (riemannian_distances(log_fwd) + riemannian_distances(log_bwd)) / 2.0
\end{verbatim}
i.e., the average of forward ($z_i$-based) and backward ($z_j$-based) linearized
distances. This is a sensible symmetric approximation (both terms have error
$O(\Kmax r^3)$ by Lemma 1) but is not documented in the paper. Specifically,
the symmetrized distance satisfies:
\[
    \frac{w_{ij}^{fwd} + w_{ij}^{bwd}}{2}
    = \dR^*(z_i,z_j) + O(\Kmax r^3)
    \quad\text{(Lemma 1 applied to both)}
\]
so the error bound is unchanged. The paper should acknowledge this in the
E-step description.
\end{bugbox}

\begin{bugbox}[B4: Candidate set for E-step is not exhaustive by default]
The trainer's default is \texttt{n\_extra\_candidates = 0}, meaning only edges
already in the current graph are reconsidered during the E-step. Once an edge
$(i,j)$ is dropped from the graph (e.g., due to a large linearization artifact),
it cannot be re-added in subsequent iterations. This violates the graph-recovery
property assumed in the convergence analysis (which requires the E-step to
consider all potentially neighboring pairs). The paper's Step 2 of the algorithm
says ``current edges plus random candidates,'' but the default setting uses no
random candidates. For the torus, where 2-cycles can occur, this means a
dropped edge stays dropped permanently even after the decoder has improved.
\end{bugbox}

\begin{bugbox}[B5: Torus geodesic evaluation uses flat formula on round data]
\texttt{compute\_true\_geodesic\_distances} for \texttt{"torus"} returns:
\begin{verbatim}
D = np.sqrt((R_torus * dtheta)**2 + (r_torus * dphi)**2)
\end{verbatim}
This is the geodesic distance on the \emph{flat} torus $(\R/2\pi\Z)^2$ with
metric $(R\,d\theta)^2 + (r\,d\phi)^2$. However, the data is generated on the
\emph{embedded} torus $(R+r\cos\phi)\cos\theta, \ldots$ in $\R^3$ via the
standard parametric embedding, which has non-flat metric
$ds^2 = (R+r\cos\phi)^2 d\theta^2 + r^2 d\phi^2$ (not flat; not the formula
used). The isometry MAE reported in Table 1 for the torus compares the decoder's
Riemannian distances to the wrong ground truth.
\end{bugbox}

\subsection{Theoretical Inconsistency in Code Comments}

\begin{bugbox}[B6: Code comment claims $\hat\kappa = 2\sqrt{K}$ for all manifolds]
\texttt{curvature.py} line 182:
\begin{verbatim}
# kappa_hat = 2*sqrt(K) for constant-curvature manifolds (sphere, flat torus).
\end{verbatim}
As proved in Proposition~\ref{prop:kappa_scale}, this holds only when the decoder
is isometrically calibrated. For an uncalibrated decoder with scale $\alpha \ne
\sqrt{3/G}$, the value is $(2/R)\sqrt{3/G}\cdot(\sqrt{3/G}/\alpha) \ne 2\sqrt{K}$.
The comment conflates two regimes. Additionally, ``flat torus'' has $K = 0$, so
$2\sqrt{K} = 0$, but the embedded torus used in experiments has non-zero local $K$.
\end{bugbox}

%% ====================================================================
\section{Sphere vs.\ Torus: Deep Mathematical Analysis}
%% ====================================================================

The paper's Section 5.5 gives three practical reasons for the sphere outperforming
the torus. Below we provide the deeper mathematical explanation, organized by
source of difficulty.

\subsection{Topological Obstruction: $\pi_1(T^2) = \Z \times \Z$}

The sphere $S^2$ is simply connected: $\pi_1(S^2) = 0$. Every loop on $S^2$
is contractible. A 2D latent space $\R^2$ with the Gaussian prior $\mathcal{N}(0,I)$
can represent $S^2$ via a two-chart atlas (north pole neighborhood + south pole
neighborhood), and the KL regularization naturally compresses the representation
toward a topological disk -- compatible with the local structure of $S^2$.

The flat torus $T^2$ has fundamental group $\pi_1(T^2) = \Z \times \Z$
(generated by two non-contractible loops $\gamma_1, \gamma_2$). Any continuous
map $\R^2 \to T^2$ has a well-defined ``winding number'' for each generator.
A VAE encoder $g_\phi: \R^G \to \R^2$ must ``cut open'' the torus along these
two loops to map it into the plane, introducing artificial boundary discontinuities.

\paragraph{Consequence.}
Points near the two cut boundaries are genuinely close on $T^2$ (geodesic distance
$\approx 0$ across the boundary) but receive latent codes far apart in $\R^2$
(latent distance $\approx L$, the period of the encoding). The Riemannian loss
penalizes $\|W e_{ij} - \ell_{ij}\|^2$, but $\ell_{ij}$ for these cross-boundary
pairs is approximately zero (nearby on manifold $\to$ small log map), while
$e_{ij} = h_\psi(\mu_j - \mu_i)$ encodes the large latent difference. This
creates an irreconcilable contradiction in the Riemannian loss that is absent
for the sphere.

\paragraph{Mathematical fix.}
Use a latent space with the correct topology: $S^1 \times S^1 \subset \R^4$,
parameterized by $(\cos\theta, \sin\theta, \cos\phi, \sin\phi)$. This requires:
(i) a 4D latent space; (ii) a von Mises or wrapped Gaussian prior on each $S^1$
factor; (iii) a modified edge encoder that computes angular differences modulo
$2\pi$. Under this architecture, the torus topology is exactly represented and
the cross-boundary issue disappears.

\subsection{Wrong Manifold: Round vs.\ Flat Torus}

\paragraph{The fundamental mismatch.}
The data in the experiment is generated on the \emph{round} (embedded) torus
with Gaussian curvature:
\[
    K(\theta, \phi) = \frac{\cos\phi}{r(R + r\cos\phi)},
    \quad R = 2, r = 1.
\]
At $\phi = 0$ (outer equator): $K_{\max} = 1/(1 \cdot (2+1)) = 1/3 > 0$.
At $\phi = \pi$ (inner equator): $K_{\min} = -1/(1 \cdot (2-1)) = -1 < 0$.

The theoretical prediction $\hat\kappa^* = 0$ (flat) is wrong for this data.
The true expected curvature proxy (averaged over triangles uniformly distributed
on the embedded torus) is a positive weighted average of $|K(\theta,\phi)|$:
\[
    \E[\hat\kappa] \;\approx\; C \int_0^{2\pi}\int_0^{2\pi}
    \frac{|\cos\phi|}{r(R+r\cos\phi)} \cdot (R + r\cos\phi) \,r\, d\theta\, d\phi
    \;=\; C' \cdot 2r \int_0^{2\pi} |\cos\phi|\, d\phi = C' \cdot 8r > 0.
\]
The torus experiment cannot converge to $\hat\kappa = 0$ because the data
manifold is not flat. The large baseline value ($\hat\kappa = 2.554$) partly
reflects this genuine non-flatness of the round torus.

\paragraph{Mathematical fix.}
Generate data on the \emph{flat} torus by first parameterizing
$(\theta, \phi) \in [0,2\pi)^2$ uniformly, then embedding via
$(\cos\theta, \sin\theta, \cos\phi, \sin\phi)/\sqrt{2} \in S^1 \times S^1 \subset \R^4$
(which has flat induced metric), then projecting to $\R^G$ via a random matrix
$A \in \R^{G \times 4}$. This yields a genuinely flat torus with $K = 0$ everywhere,
making the theoretical prediction $\hat\kappa^* = 0$ valid.

\subsection{Linearization Artifact Scale: $O(K r^3)$ vs.\ Torus Topology}

For the sphere with $K = 1/R^2 = 1$ and KNN radius $r_n$, the linearization error
per edge is $O(r_n^3)$. For the round torus, $K$ varies in $[-1, 1/3]$, and the
encoder's ``cut-open'' representation means the effective KNN radius in latent
space is not the manifold geodesic radius. Specifically, for points near the
cut boundary, $\|\Delta z_{ij}\|$ can be $O(L)$ (full period length) while
$d_M(x_i, x_j) = O(\varepsilon)$. In this case the linearization error is
$O(\|\Delta z\|^3) = O(L^3)$, orders of magnitude larger than $O(\varepsilon^3)$.

This explains the 2-cycle behavior observed in the torus experiment: the KNN
graph oscillates between two states because the Riemannian distances for
cross-boundary pairs are dominated by linearization artifacts, not by true
Riemannian geometry.

\paragraph{Formal bound.}
Let $L = O(\sqrt{Rr} \cdot 2\pi)$ be the typical cross-boundary latent distance.
By Lemma 1, the linearization error for a cross-boundary pair is:
\[
    |\tilde{w}_{ij} - \dR^*(z_i,z_j)| \le \frac{\Kmax}{6} L^3 \gg \frac{\Kmax}{6} r_n^3.
\]
The distance-clipping at $5 \times \mathrm{median}(w_{ij})$ prevents the worst
cross-boundary pairs from entering the KNN but does not resolve the underlying
topological mismatch.

\subsection{Holonomy of Non-Contractible Loops}

The curvature proxy $c_{ijk}$ measures holonomy around \emph{contractible}
triangles. For the flat torus, holonomy around any contractible loop is zero
($K = 0$). However, the two generators $\gamma_1, \gamma_2$ of $\pi_1(T^2)$
carry non-trivial information: going around each generator and returning encodes
the periodic structure. The self-consistent iteration, which is based on
triangle-level terms, has no mechanism to detect or exploit this non-contractible
holonomy. The sphere, being simply connected, has no such non-contractible loops,
and all its geometric information is encoded in local triangle terms.

\paragraph{Mathematical consequence.}
Even at the fixed point $(f_\theta^*, G^*)$, the SCR-VAE's curvature proxy
provides no signal about the torus's topological structure. The edge decoder $W$
spans the local tangent frame but cannot represent the global topology. This is
why the torus requires many more iterations ($\approx 98$ vs. $\approx 15$
for the sphere): the algorithm must resolve the global topology via many local
corrections, rather than exploiting a topological shortcut.

\subsection{KL Prior vs.\ Torus Geometry}

The Gaussian prior $\mathcal{N}(0, I_d)$ in the node encoder's ELBO pulls all
latent codes toward the origin. For the sphere (which can be represented as a
punctured plane via stereographic projection, with the removed point corresponding
to the north pole pushed to infinity), this is roughly compatible: most data
points map to a bounded region of the latent plane.

For the torus, the Gaussian prior creates a systematic \emph{asymmetry} between
the two periodic directions: the direction pointing toward the origin (radially)
gets compressed by the KL penalty, while the orthogonal direction is less
penalized. This breaks the rotational symmetry of the torus and causes one
periodic direction to be learned less well than the other, manifesting as
different KNN behavior along the two periodic axes.

\paragraph{Quantitative estimate.}
Let the torus be encoded as a rectangle $[0, L_1] \times [0, L_2]$ in latent
space. The KL term contributes $\sim L_1^2 + L_2^2$ to the loss (for Gaussian
prior). The optimal encoding under the KL constraint sets $L_1 = L_2 = L$
(square), but the optimal encoding under the Riemannian loss may prefer $L_1/L_2
= R/r = 2$ (matching the aspect ratio of the torus). These two objectives
conflict, and the equilibrium is reached at some intermediate aspect ratio that
satisfies neither constraint exactly.

\subsection{Summary: Why the Sphere Succeeds Where the Torus Struggles}

\begin{table}[h]
\centering
\caption{Root causes of the sphere vs.\ torus performance asymmetry.}
\begin{tabular}{lcc}
\toprule
Source of Difficulty & Sphere $S^2$ & Torus $T^2$ \\
\midrule
Fundamental group $\pi_1$ & $0$ (trivial) & $\Z \times \Z$ (complex) \\
Latent representation with $\R^2$ encoder & Compatible & Requires ``cut'' \\
Cross-boundary artifact pairs & None & Severe \\
Ground truth curvature $K$ & $1$ (uniform) & $\cos\phi / [r(R+r\cos\phi)]$ (varying) \\
Experiment matches theory & Yes (round $S^2$ = true $S^2$) & No (round $\ne$ flat) \\
True $\hat\kappa$ target & $2/R$ (achievable) & Not $0$ (wrong target) \\
Iterations to fixed point & $\approx 15$ & $\approx 98$ \\
KL prior compatibility & Good & Asymmetric \\
\bottomrule
\end{tabular}
\end{table}

%% ====================================================================
\section{Recommended Action Items}
%% ====================================================================

\subsection{Theory Revisions (Required)}

\begin{enumerate}
\item \textbf{Add Proposition~\ref{prop:existence}} (fixed-point existence) to
Section 4, before the Fixed-Point Characterization Theorem.

\item \textbf{Split the PL assumption} into (PL) and (SC) as described in the
Gap 3 bridge. Clarify that $L_*$ follows from (SC), not from (PL).

\item \textbf{Add Proposition~\ref{prop:a4prime}} (deriving A4' from
$\delta_{\mathrm{rec}} = O(r_n^2)$) as a remark or lemma following Assumption A4'.

\item \textbf{Fix Equation (23):} Replace $8R^2 K A$ with $4\sqrt{K} A$.

\item \textbf{Fix Section 5.5 embedding claim:} Replace ``cannot be embedded in
any finite dimension'' with ``cannot be smoothly embedded in $\R^3$.''

\item \textbf{Reconcile curvature proxy formulae:} Add
Proposition~\ref{prop:kappa_scale} (or a remark) explaining that $\hat\kappa
\to 2/R$ at the isometric fixed point and $(2/R)\sqrt{3/G}$ for uncalibrated
decoders. Clarify that the ``True'' column in Table 1 is the isometric target.

\item \textbf{Add the edge KL term to Equation (3)} and briefly analyze its
role (variational regularization of $h_\psi$).

\item \textbf{Fix the torus experiment description:} Acknowledge that the data
is on the round (not flat) torus, and that the flat-torus geodesic formula is
an approximation valid for small $\Delta\theta, \Delta\phi$.
\end{enumerate}

\subsection{Theory Extensions (Suggested)}

\begin{enumerate}
\item \textbf{Open Problem resolution for torus:} Propose the $S^1 \times S^1
\subset \R^4$ latent parameterization as the correct fix. Derive the corresponding
von Mises prior and show that the self-consistent iteration for this architecture
converges without the cross-boundary artifact.

\item \textbf{Non-contractible holonomy as an additional curvature certificate:}
For manifolds with non-trivial $\pi_1$, the closure of the decoder around
non-contractible loops $\gamma$ provides topological information:
$\sum_{(i,j) \in \gamma} \ell_{ij}$ is the holonomy of $\gamma$ in the decoder
manifold, which equals zero for a flat connection. This could extend the curvature
certificate to detect topology, not just curvature.

\item \textbf{Lower bounds on isometry error:} Under what conditions is it
impossible to achieve $\varepsilon$-isometry with latent dimension $d$? For the
torus in $\R^2$ latent space, prove a lower bound on the isometry error arising
purely from the topological mismatch.
\end{enumerate}

\subsection{Code Fixes (Required Before Release)}

\begin{enumerate}
\item \textbf{Fix the broken test} \texttt{test\_edge\_decoder\_no\_bias}: replace
\texttt{edge\_decoder.W.bias} with a correct existence check on the edge decoder's
parameters.

\item \textbf{Document symmetrized E-step distance}: Add a comment or footnote
acknowledging the symmetrization $(\dR_{fwd} + \dR_{bwd})/2$ and note that both
terms have error $O(\Kmax r^3)$ by Lemma 1.

\item \textbf{Default \texttt{n\_extra\_candidates > 0}}: Set
\texttt{n\_extra\_candidates = 50} (or proportional to $N$) by default to allow
the graph to discover dropped edges in subsequent iterations.

\item \textbf{Fix torus geodesic formula} in \texttt{compute\_true\_geodesic\_distances}:
either use the correct embedded torus metric
$ds^2 = (R + r\cos\phi)^2 d\theta^2 + r^2 d\phi^2$ (integrated numerically),
or clearly document that the flat formula is used as an approximation.

\item \textbf{Correct code comment in} \texttt{curvature.py}: Replace
``$\hat\kappa = 2\sqrt{K}$ for constant-curvature manifolds'' with ``$\hat\kappa
\to 2\sqrt{K}$ at the isometrically-calibrated fixed point; value depends on
decoder scale otherwise.''
\end{enumerate}

%% ====================================================================
\section{Proof Status Summary}
%% ====================================================================

\begin{center}
\begin{tabular}{lll}
\toprule
Statement & Status & Key condition / issue \\
\midrule
Lemma 1: linearized distance error $\le K_{\max}r^3/6$
  & \textbf{Correct} & Bertrand-Puiseux in normal coords \\
Proposition 1: $\col(W^*) = \mathrm{top}$-$k(M_\theta)$
  & \textbf{Correct} & Eckart-Young \\
Proposition 2 (Ky Fan)
  & \textbf{Correct} & Standard \\
Theorem 1(a): $w^*_{ij} = \|\ell^*_{ij}\|$
  & \textbf{Correct} & Definition \\
Theorem 1(b): $|w^*_{ij} - \dR^*| \le K_{\max}r^3/6$
  & \textbf{Correct} & Lemma 1 \\
Theorem 1(c): $\col(W^*)$ correct
  & \textbf{Correct} & Prop.\ 1 \\
\textbf{Existence of fixed point}
  & \textbf{GAP -- bridge provided} & See Prop.\ \ref{prop:existence} \\
Proposition 3 convergence: $\Delta^* = O(r_n^2)$
  & \textbf{Conditional} & PL $+$ SC (circular; see Gap 3) \\
Theorem 2 isometry bound
  & \textbf{Conditional on A4'} & See Gap 2, Prop.\ \ref{prop:a4prime} \\
Proposition 4 (closure vector)
  & \textbf{Correct} & Hessian Taylor, Schwarz \\
Proposition 5 (curvature detection)
  & \textbf{Correct with caveat} & $K=0 \Leftrightarrow$ locally affine, not $\Leftrightarrow K = 0$ \\
Curvature proxy $\hat\kappa \to 2\sqrt{K}$
  & \textbf{Correct only at isometric fixed pt} & See Prop.\ \ref{prop:kappa_scale} \\
Eq.\ (23) coefficient $8R^2KA$
  & \textbf{WRONG -- should be $4\sqrt{K}A$} & Sign/factor error \\
Flat torus embedding claim
  & \textbf{WRONG -- correct is $\R^3$ only} & $\R^4$ isometric embedding exists \\
\bottomrule
\end{tabular}
\end{center}

%% ====================================================================
\section{Conclusion}
%% ====================================================================

The SCR-VAE paper makes a genuine and significant contribution: the self-consistent
co-evolution of the KNN graph and the decoder's Riemannian geometry is a principled
and novel approach to isometric generative modeling. The core theoretical results
(Lemma 1, Proposition 1, Theorem 1, Proposition 4) are correct and non-trivial.

The identified gaps -- most critically the absence of a fixed-point existence proof,
the circular structure of the convergence proposition, the curvature proxy formula
inconsistency, and the fundamental mismatch between the experimental torus setup
and the theoretical claims -- require revision before publication.

The sphere vs.\ torus performance gap is deeper than the paper acknowledges. The
primary causes are: (1) a topological obstruction from $\pi_1(T^2) = \Z \times \Z$
that forces artificial cut-boundary artifacts in the 2D latent representation;
(2) the use of the wrong torus model (round vs.\ flat) in both data generation
and evaluation; and (3) the incompatibility of the Gaussian prior with the torus
topology. The mathematical fix -- using a $S^1 \times S^1 \subset \R^4$ latent
parameterization with a von Mises prior -- would likely eliminate most of the
performance gap and strengthen the paper's empirical claims.

With the identified corrections and the suggested extensions, this paper would
represent a strong contribution to the NeurIPS 2026 program.

\end{document}
